\scanpage{158}
\mspage{107}

Faisons un énoncé en forme

\uline{Proposition 11} On peut trouver une situation \unsure{$(x)$|$(\alpha)$} (p. 124) avec

a) $t$ objet final de $X$

b) $U,U'$ (et $X$) $\unsure{W|\mathcal W}_\infty$-asphériques, donc $i,p\in W_\infty$ \inserted{(et même $i$ $W_{\unsure{\infty|0}}$-asphérique)}

c) $X'$ non $W_\infty$-asphérique, donc $i',q\notin W_\infty$.

Il suffit pour cela de partir d'un objet 1-connexe \deleted{et non $\infty$-connexe} $Y$ de $\mathrm{Cat}$, de lui adjoindre un objet final \inserted{$u'_0$}, ce qui donne $U'$, de prendre $U=\underline{\mathrm{Ch}}(Y;y_0,-)$, $p:U\to U'$ le composé $U\xrightarrow{\mathrm{but}}Y\mathrel{\mathop{\hookrightarrow}\limits^{\mathrm{incl.}}_{\mathrm{can}}}U'$, et $X$ \unsure{se} déduit de $U$ en lui ajoutant un objet final $t$. On a alors $U'_{/t'}\simeq Y$, \deleted{(\unsure{donc} si $Y$ était $\infty$-connexe)} d'où

\[
 X'\ W_\infty\text{-asph}\quad (\text{i.e. }i'\in W_\infty)\Longleftrightarrow Y\quad W_\infty\text{-asphérique.}
\]

On a le contr'exemple cherché en prenant $Y$ \deleted{(1-connexe mais)} non $\infty$-connexe.

\uline{2.4.} \uline{Application aux immersions ouvertes dwyériennes}

Donnons avant $=$ l'énoncé \uline{positif} auquel nous en \unsure{arrivons|étions}.

\uline{Théorème 3} Soient $(X,U,T)$ un triple $\in\unsure{\mathscr T|\mathcal T}$ ($U$ ouvert de $X$ dans $\mathrm{Cat}$, $T=X-U$), \inserted{et pour \unsure{la} \deleted{\unsure{toute}} flèche} $p:U\longrightarrow U'$ \deleted{\unsure{une}} flèche de $\mathrm{Cat}$, \deleted{\unsure{soit}} le carré cocartésien dans $\mathrm{Cat}$

\diagram{\begin{tikzcd}[ampersand replacement=\&,column sep=small]
U \arrow[r,hook,"i"] \arrow[d,"p"'] \& X \arrow[d,"q"'] \& T \arrow[l,dashed,hook] \arrow[d,dashed,"\mathrm{iso}"]\\
U' \arrow[r,"i'"] \& X' \& T' \arrow[l,dashed,hook]
\end{tikzcd}}
% archive sheet number: 158
